% Transcription — fonds Alexandre Grothendieck, shelfmark 69, batch 4 (pages 61-80).
% Established from the facsimile archives/batches/69/batch-04.pdf (cut from a
% copy of 69.pdf fetched through the project relay, a mirror of
% https://grothendieck.umontpellier.fr/69.pdf), per the skill
% transcribe-grothendieck. The folder is 119 pages in six batches (1-20, 21-40,
% 41-60, 61-80, 81-100, 101-119), transcribed in parallel; this is the fourth.
% The raw scan's cover sheet is not in the batch file: PDF page k is
% archivists' page 60+k, and the pencilled numbers bottom-left (« 61 » ... « 80 »)
% agree.
%
% Pass: Opus 5.5 (claude-opus-5-5), 2026-09-27 — first pass, unchecked against the pages by a human.
%
% Sixteen of the twenty pages carry a \page marker. Absent: page 64, a pink
% separator sheet; pages 65, 72 and 76, cream separator sheets bearing only
% show-through of the neighbouring leaves. Nothing typed, no letter, nothing
% by another hand in this batch. Pages summarised: none; pages 77-80 are
% sheets of figures and scattered formulas, transcribed as far as they read,
% with the figures described in \drawing and the unreadable fragments noted.
%
% His pagination: none on these leaves. Page 66 opens with his underlined
% heading « Parties i-libres etc. », made the \section; pages 66-71 run on
% continuously (« (cf. plus bas) », « il fallait commencer par là ! »), 73-75
% form a second run (bitriangles), 75 breaking off after two lines. No dates.
% Paper: 61 is punched computer-listing paper; 62-63 are half-sheets of USTL
% « Institut de Mathématiques » letterhead (printed heading only, used as
% scrap; 62 is written in both directions); the rest plain sheets.
%
% What the batch is about. The combinatorics of the graph C of the 27 lines
% on a cubic surface (vertices « liés » when the lines meet: each vertex has
% 10 neighbours, 16 non-neighbours), with W = Aut(C) of order 51 840 (he
% writes W_{SO(10)}, 2^7 3^4 5), and the root system E_6 behind it.
%   61     A table of « types » of configurations in C (1- to 6-bases,
%          ordered or not, half-stars, triads, triangles, arcs): for each, the
%          equivalent category (Ens_5 x Ens_2 ...), its automorphism group, the
%          number of such configurations (27, 216, 432, 720, 4320, 1080,
%          25 920, 51 840, 72 ...) and the orbit lengths.
%   62-63  The roots in the basis eta, xi_1..xi_6 (xi = sum): r_0 = 2 eta - xi,
%          r_ij = xi_i - xi_j, r_ijk = eta - xi_i - xi_j - xi_k, the double-sixes
%          attached to them, the table of scalar products r.r' in {0, +-1}, the
%          A_2 stars r_0, r_123 = r_6, r_456, and r_0 = 3(r_1+r_2+r_3)+2r_4+r_5+2r_6.
%   66-71  « Parties i-libres »: n-free subsets (pairwise non-linked) exist
%          iff n <= 6; maximal ones are the 6-bases and the 2-stars
%          E(x) ∩ E(y); the tower Pl(L_6,C) -> ... -> Pl(L_1,C) = C with degrees
%          27, 16, 10, 6, 2, 1 and the resulting counts B_i(C); case by case,
%          pairs (via bibases {b,b'} and xi_ij), 3-free sets (bijection with
%          couples (B,K), 720 of them, equivalence with triples (tau,K,K'),
%          involution K <-> K', J exactly in two bases, pairs of roots at 60°),
%          hexagons L(J), 4-free and 5-free sets, 2-stars; then « Bases et
%          bibases (il fallait commencer par là !) »: 72 bases, 36 bibases.
%   73-75  Bitriangles Delta = S x T (9 vertices, 18 edges) and the « bloc de
%          triades »; Proposition: the complement of a bitriangle is a
%          bitriangle, proved through the six bijections S -> T split by parity
%          into two families; conclusion: Delta -> Delta-bar is an involutive
%          autoequivalence of (Bitr), i.e. of sets of 6 with a (3,3) partition;
%          a « duality » between T and S, left unfinished on 75.
%   77-80  Sheets of drawings and counts: polyhedra in the cubic complex
%          (prisms, antiprisms, pseudocubes, cubes 27.10.2, skew dodecahedra
%          2^4 3^3), polygons (squares, pentagons, 720 hexagons), a labelled
%          skew dodecahedron; upside down on 77, Z[C] -> Ê(C), the form
%          c.c' in {-1,0,1} and E(C) = lambda^perp; on 78, M ⊂ F_2^Q and
%          e_i-computations; on 79, 45.8.. = 360, an order-144 automorphism
%          group, « pentaprismes » as couples of 3-bases, and U,V,W with
%          UVW in {1, sigma}; on 80, a « Lemme » on triangles u,v,w, left as a
%          question (an implication struck through).
%
% Notation, fixed once. C the cubic graph (« complexe cubique »), vertices
% « liés » / « non liés »; E(x) the neighbours of x; xi_i, xi'_i, xi_ij his
% labels of the 27 vertices (a double-six and the 15 others), eta the class
% he adjoins on 62-63; b, b' bases, B = {b,b'} a bibase, sigma_B, I = B/sigma_B;
% Pl(L_i,C) the embeddings of the discrete graph L_i, Pl' and Pl'' the two
% orbits for i = 5; B_i(C), Bib(C); L(J) as \mathcal{L}; Ens_n the groupoid of
% n-sets; \mathfrak{S}_n for his gothic S; \mathfrak{P}_2 for the set of
% 2-subsets; tau, K, K', S, T as written. The table of page 61 is set as a
% list, the columns separated by vertical bars, since it does not fit a page
% measure as a tabular.
%
% Apparatus: 69 \ill{}, 45 \uncertain{} (counted on the final file). Hardest
% pages: 62 (two directions, overwritten coefficients), 63 (the margin
% column, unread), 70 (dense, two columns, a boxed and partly struck margin
% note), 77-80 (fragments among drawings). Readings to check first: « complexes
% cubiques » on 66; the struck line and the last sentence of 69; the E(J)
% sentence on 70; « quasi-prédéterminée » on 74; the i ≠ i' of 63.
%
% The dating is the folder's own.
\documentclass[11pt,a4paper]{article}
\input{../preamble/grothendieck}

\folder{69}
\batch{4}
\pages{61}{80}
\dating{[à partir de 1976]}
\watermark{Édition de démonstration}
\foldertitle{Graphes cubiques : notes manuscrites (s.d.), lettre (s.d.).}

\begin{document}

\page{61}
\note{page entière occupée par un tableau à cinq colonnes, lignes numérotées de $1^\circ$ à $12^\circ$ dans la marge gauche, puis quelques lignes non numérotées. En-têtes : « Cat.\ à étudier des » \,|\, « équivalente à » \,|\, « groupe d'automorphismes » \,|\, « cardinal de l'ens.\ des él.\ de structure du type envisagé » \,|\, (sans titre : des suites d'entiers, sans doute les longueurs des orbites). On donne le tableau ligne par ligne, les colonnes séparées par des barres verticales. Chaque type est accompagné d'un petit schéma de points (un, deux, \ldots, six points, parfois numérotés), qu'on ne reproduit pas.}

\begin{itemize}
\item[$1^\circ$] $1$-base \,|\, \struck{\ill{}} ens.\ $J$ de card.\ $5$ \uncertain{et} torseur sous $\mathbb{F}_2(J)'$ \,|\, $W_{SO(10)}$, ext.\ de $\mathfrak{S}_5$ par $(\mathbb{Z}/2\mathbb{Z})^4$, de cardinal $2^4\,5!=2^7\cdot3\cdot5=1920$ \,|\, $3^3=27$ \,|\, $1,10,16$
\item[$2^\circ$] $2$-base \,|\, $\mathrm{Ens}_5\times\mathrm{Ens}_2$ \,|\, $\mathfrak{S}_5\times\mathfrak{S}_2$, de card.\ $2\cdot5!=2^4\cdot3\cdot5=240$ \,|\, $2^3 3^3=27\cdot8=216$ \,|\, $1,5,10,10$
\item[$3^\circ$] $2$ bases ordonnées \,|\, $\mathrm{Ens}_5$ \,|\, $\mathfrak{S}_5$ de card.\ $5!=120$ \,|\, $2^4 3^3=27\cdot16=432=6\cdot72$ \,|\, $1,5,5,5,10$
\item[$4^\circ$] $3$ base \,|\, $\mathrm{Ens}_3\times\mathrm{Ens}_3\times\mathrm{Ens}_2$ \,|\, $\mathfrak{S}_3\times\mathfrak{S}_3\times\mathfrak{S}_2$, de card.\ $6\cdot6\cdot2=72=2^3 3^2$ \,|\, $2^4 3^2 5=16\cdot45=720$ \,|\, $3,3,6,6,9$
\item[$5^\circ$] $3$ bases ordonnées \,|\, $\mathrm{Ens}_3\times\mathrm{Ens}_2$ \,|\, $\mathfrak{S}_3\times\mathfrak{S}_2$, card.\ $12$ \,|\, $2^5 3^3\cdot5=27\cdot16\cdot10=4320$ \,|\, \uncertain{$1,1,1,2,2,2,2,3,3,3,3,6$}
\item[$6^\circ$] $4$-base \,|\, $\mathrm{Ens}_4\times\mathrm{Ens}_2$ \,|\, $\mathfrak{S}_4\times\mathfrak{S}_2$, card.\ $2\cdot4!=48=2^4 3$ \,|\, $2^3 3^3 5=27\cdot16\cdot10\cdot6=1080$ \,|\, \uncertain{$1,2,3,4,4,6,8$}
\item[$7^\circ$] $4$ bases ordonnées \,|\, $\mathrm{Ens}_2$ \,|\, $\mathfrak{S}_2$, card.\ $2$ \,|\, $2^6\cdot3^4\cdot5=25\,920$ \,|\, $15\times1$, $6\times2$
\item[$8^\circ$] $5$ base \,|\, $\mathrm{Ens}_5$ \,|\, $\mathfrak{S}_5$, card.\ $5!=120=2^3\cdot3\cdot5$ \,|\, $2^4\cdot3^3=27\cdot16=432=6\cdot72$ \,|\, $1,1,5,5,5,10$
\item[$9^\circ$] $5$ bases ordonnées \,|\, $e$ \,|\, $e$, card.\ $1$ \,|\, $2^7\cdot3^4\cdot5=51\,840$ \,|\, $27\times1$
\item[$10^\circ$] $6$ bases \,|\, $\mathrm{Ens}_6$ \,|\, $\mathfrak{S}_6$, card.\ $720=2^4 3^2 5$ \,|\, $2^3 3^2=8\cdot9=72$ \,|\, $6,6,15$
\item[$11^\circ$] $6$ bases ordonnées \,|\, $e$ \,|\, $e$, card.\ $1$ \,|\, $2^7\cdot3^4\cdot5=51\,840$ \,|\, $27\times1$
\item[$12^\circ$] $\frac12$ étoile, formée \uncertain{de} $5$-bases \,|\, $\mathrm{Ens}_5\times\mathrm{Ens}_2$ \,|\, cf.\ plus haut $2^\circ$ \,|\, \uncertain{$15$} \,|\, $1,1,5,5,5,10$
\item[$*$] triades ordonnées \,|\, \struck{\ill{}} \,|\, $\mathfrak{S}_3\times\mathfrak{S}_3\times\mathfrak{S}_3$, ordre $6^3=2^3 3^3=216$ \,|\, $2^4\cdot3\cdot5=240$ \,|\, $9,9,9$
\item[] triades \,|\, perm.\ d'ordre $3$ d'ens.\ de card.\ $3$ ; ens.\ de card.\ $9$ avec partition \,|\, $\mathfrak{S}_3\cdot(\mathfrak{S}_3\times\mathfrak{S}_3\times\mathfrak{S}_3)$, card.\ $6^4=2^4 3^4$ \,|\, $2^3\cdot5=40$ \,|\, $27$
\item[] grilles $3$-carrées \,|\, (somme directe avec partition \uncertain{distinguée}) \,|\, $\mathfrak{S}_2\cdot(\mathfrak{S}_3\times\mathfrak{S}_3\times\mathfrak{S}_3)$ \,|\, $2^3\cdot3\cdot5=120$ \,|\, $9,18$
\item[] arête $\sim$ \uncertain{triangles pointés} \item[] arcs $\to$ $\sim$ triangles ordonnés \item[] triangles $\Delta$ \item[] $V$ \item[] $(\bullet\!\to\!\bullet)\simeq$(arcs)
\end{itemize}
\note{la ligne $12^\circ$ porte un petit dessin (une demi-étoile : un losange avec des arcs issus d'un sommet). La ligne marquée d'un astérisque est celle des triades ordonnées ; en marge, en face de la dernière ligne, « $13^\circ$ » (lecture douteuse) et « $16$ ». Les suites de la dernière colonne des lignes $5^\circ$ et $6^\circ$ sont écrites sur deux lignes serrées. Dans la ligne $8^\circ$, $27\cdot16=432$ est surchargé. Dans la ligne des triades, « $2^3\cdot5$ » et « $6^4$ » sont récrits.}

\page{62}
\note{brouillon de calculs sur une demi-feuille à en-tête imprimé, écrit dans les deux sens et traversé de deux traits obliques ; on donne d'abord le côté à l'endroit, puis le côté tête-bêche.}

$\delta_1-\delta_8\,\delta_9$ ; $\delta_1-\delta_n\,\delta_{n+1}$, $n+1+g\leq27$, $n\leq\ill{}$
\[
r_1+r_2+r_3,\qquad r_2+r_3+r_4,\qquad r_3+r_4+r_5 ;
\]
\struck{$(r_1+r_{i-1})+$} $(r_1+\cdots+r_{i-1})+(r_2+\cdots+r_{j-1})+(r_3+\cdots+r_{k-1})$ ($i\geq3$), $r_1+2(r_2+\cdots+r_{i-1})$, $r_1+2r_2$ ; $i-1$, $j-2$, $k-3$ ; $i+j+k-6$.

\drawing{trois petits graphes en zigzag (lettres « M » et « W » faites de sommets et d'arêtes), deux d'entre eux entourés de boucles, une ligne brisée en créneau, un chevron.}

Tête-bêche :
\[
\begin{array}{ll}
r_0=2\eta-\ldots & \\
r_1=\xi_1-\xi_2 & \\
r_2=\xi_2-\xi_3 & 3(r_1+r_2+r_3)=3(\xi_1-\xi_4)\\
r_3=\xi_3-\xi_4 & \\
r_4=\xi_4-\xi_5 & 2r_4=\ldots\xi_4-2\xi_5\\
r_5=\xi_5-\xi_6 & 2r_5=\ldots\xi_5-\xi_6\\
r_6=\eta-\xi_1-\xi_2-\xi_3 &
\end{array}
\qquad 3\xi_1-\xi_4-\xi_5-\xi_6
\]
\note{les coefficients marqués « $\ldots$ » sont surchargés et ne se lisent pas ; avant $r_6$, une ligne biffée. Plus bas, de ce côté, des inégalités $t_1\geq\ldots$, \uncertain{$t_2\geq6+1+n$}, et deux petits graphes entourés, aux sommets marqués $\xi$.}

\page{63}
\[
r_0=2\eta-\xi,\qquad r_{ij}=\xi_i-\xi_j\quad(i<j),
\]
$r_{\lbrace i,j,k\rbrace}=$ \struck{$(\xi_i,\xi_j,\xi_k,\xi_{mn},\xi_{n\ell},\xi_{\ell m})$}
\[
r_{\lbrace i,j,k\rbrace}=\xi_{jk}-\xi_i=\xi_{ki}-\xi_j=\xi_{ij}-\xi_k
\]
\[
=\xi'_\ell-\xi_{mn}=\xi'_m-\xi_{n\ell}=\xi'_n-\xi_{\ell m}=\eta-\xi_i-\xi_j-\xi_k
\]
\note{sous la parenthèse biffée, une seconde ligne également biffée : $\xi_{jk},\xi_{ki},\xi_{ij},\xi'_\ell,\xi'_m,\xi'_n$. En face, à droite, les couples « $\binom{\xi_i\ \ -\xi_6}{\xi'_i\ \ -\xi'_6}$ » (lecture douteuse), puis
$\begin{pmatrix}\xi_i,\xi'_j,\xi_{j\ell}\ \ \ell\neq i,j\\ \xi_j,\xi'_i,\xi_{i\ell},\ \ell\neq i,j\end{pmatrix}$ et
$\begin{pmatrix}\xi_i,\xi_j,\xi_k, & \xi_{mn},\xi_{n\ell},\xi_{\ell m}\\ \xi_{jk},\xi_{ki},\xi_{ij}, & \xi'_\ell,\xi'_m,\xi'_n\end{pmatrix}$ : les « doubles-six » associés aux racines. Dans la marge gauche, un empilement de formules raturées, écrites en long, illisible \ill{}.}

\[
r_0\cdot r_{ij}=0,\qquad r_0\cdot r_{ijk}=-1,
\]
\[
r_{ij}\cdot r_{i'j'}=\begin{cases}0 & \lbrace i,j\rbrace\cap\lbrace i',j'\rbrace=\emptyset\\ \pm1 & \lbrace i,j\rbrace\cap\lbrace i',j'\rbrace\neq\emptyset\end{cases}
\]
et à côté : $+1$ si \struck{$i<j=i'<j'$ ou $i'<j'=i<j$} \add{$i'=j$ ou $i=j'$}, $-1$ si $i\neq i'$ ou $j=j'$.
\note{ainsi sur la page, $i\neq i'$ ; on attendrait $i=i'$. Le signe est peut-être un $=$ mal formé.}
\[
r_{ijk}\cdot r_{pqr}=1-\operatorname{card}(\lbrace i,j,k\rbrace\cap\lbrace p,q,r\rbrace)=\begin{cases}-1 & \operatorname{card}(\lbrace i,j,k\rbrace\cap\lbrace p,q,r\rbrace)=2\\ 0 & \ldots=1\\ 1 & \ldots=0\end{cases}
\]
\[
r_{ij}\cdot r_{pqr}=\begin{cases}0 & \lbrace i,j\rbrace\subset\lbrace p,q,r\rbrace\ \text{ou}\ \lbrace i,j\rbrace\cap\lbrace p,q,r\rbrace=\emptyset\\ \pm1 & \operatorname{card}(\lbrace i,j\rbrace\cap\lbrace p,q,r\rbrace)=1\end{cases}
\qquad
\begin{cases}+1 & i\in\lbrace p,q,r\rbrace\\ -1 & j\in\lbrace p,q,r\rbrace\end{cases}
\]
\[
\frac{-\langle r,r'\rangle}{\sqrt{\langle r,r\rangle\langle r',r'\rangle}}=\frac{\langle r,r'\rangle}{-2}
\]

\begin{align*}
\xi_1\xi_2\xi_3\ \ \xi_4,\xi_5\xi_6\quad &= r_0=2\eta-\xi=3(r_1+r_2+r_3)+2r_4+r_5+2r_6\\
\xi_1\xi_2\xi_3\ \ \xi_{56}\xi_{64}\xi_{45}\quad &= r_{123}=\eta-\xi_1-\xi_2-\xi_3=r_6\\
\xi_{23},\xi_{31},\xi_{12}\ \ \xi_4\xi_5\xi_6\quad &= r_{4,5,6}=\eta-\xi_4-\xi_5-\xi_6=r_6+3(r_1+r_2+r_3)+2r_4+r_5
\end{align*}
\note{la notation « $\xi_1\xi_2\xi_3\ \ \xi_4,\xi_5\xi_6$ » etc.\ désigne sans doute la base (le double-six) qui correspond à la racine ; dans la troisième ligne, l'écriture fait voir $3(\xi_1+\xi_2+\xi_3)+2\xi_4+\xi_5$ avec des $r$ récrits sur des $\xi$.}

\drawing{deux étoiles de six rayons (système de racines de type $A_2$) : la première aux extrémités $r_{456}$, $r_0$, $r_{123}=r_6$, $-r_{4,5,6}$, $-r_0$, $-r_{123}=-r_6$, deux côtés de l'hexagone en tirets ; la seconde, en bas à droite, avec $r_{456}$, $-r_{4,5,6}$ et des indices surchargés.}

\section{Parties $i$-libres etc.}

\page{66}
\note{le titre, souligné, est de sa main, en tête de la page 66.}

On étudie des types de diagrammes plongés dans les \uncertain{complexes} cubiques $C$, liés directement aux parties $n$-libres \struck{\ill{} bases} (parties de cardinal $n$ de $C$ formées de sommets non liés deux à deux). Signalons \struck{\ill{}} de suite

\emph{Prop.} (1)\note{numéro entouré.} Pour qu'il existe en $C$ une partie $n$-libre, il faut et il suffit que $1\leq n\leq 6$. Les parties $n$-libres \uncertain{maximales} \struck{\ill{}} sont :

a) celles de cardinal $6$, ou \uncertain{sous-ensembles} des bases ;

b) celles \struck{de cardinal 5 \ill{}} \ill{} de la forme $E(x)\cap E(y)$, où $x,y$ sont deux sommets distincts non liés en $C$. Ce sont des parties de cardinal $5$.

(2)\note{numéro entouré.} Soit $L_i$, $1\leq i\leq 6$, le graphe discret $[1,i]$ \struck{de} de cardinal $i$ (pas de liens). Alors pour $i\neq 5$, l'ens.\ $\operatorname{Pl}(L_i,C)$ des plongements de $L_i$ en $C$ est un espace homogène sous $W=\operatorname{Aut}(C)$. Pour $i=5$, il a deux orbites $\operatorname{Pl}(L_5,C)'$ et $\operatorname{Pl}(L_5,C)''$ sous $W$, savoir

a) les $5$-bases, i.e.\ parties $5$-libres qui sont contenues dans une base \ill{} ;

b) les deux-étoiles, i.e.\ parties de la forme $E(x)\cap E(y)$, où $x,y$ sont deux sommets distincts non liés.

Les applications $\operatorname{Pl}$(\uncertain{homomorphismes d'inclusion})

\begin{tikzcd}[column sep=small, row sep=small, nodes={font=\scriptsize}]
 & & & & & \operatorname{Pl}''(L_5,C) \arrow[dl, "1"'] & \\
\mathrm{pt}=\operatorname{Pl}(L_0,C) & \operatorname{Pl}(L_1,C) \arrow[l, "27"'] & \operatorname{Pl}(L_2,C) \arrow[l, "16"'] & \operatorname{Pl}(L_3,C) \arrow[l, "10"'] & \operatorname{Pl}(L_4,C) \arrow[l, "6"'] & \operatorname{Pl}'(L_5,C) \arrow[l, "2"'] & \operatorname{Pl}(L_6,C) \arrow[l, "1"']
\end{tikzcd}

\note{la ligne est écrite de gauche à droite, flèches vers la gauche, les degrés au-dessus des flèches ; sous $\operatorname{Pl}(L_1,C)$ il écrit « $=$ \uncertain{$D_1$} », et « $\mathrm{pt}$ » est glosé en dessous « $=\operatorname{Pl}(L_0,C)$ ».}

sont surjectives, et les degrés de projection sont respectivement \struck{$27,16,10$} $27,16,10,6$ comme indiqué (cf.\ dessin plus bas), d'où \ill{} les cardinaux des $i$-bases \ill{}

\page{67}
\[
\begin{array}{ll}
\operatorname{Pl}(L_1,C)=C & 27=3^3\\
\uparrow\ \operatorname{Pl}(L_2,C) & 27\cdot 16=2^4 3^3=432\\
\uparrow\ \operatorname{Pl}(L_3,C) & 27\cdot16\cdot10=2^5\cdot3^3\cdot5=4320\\
\uparrow\ \operatorname{Pl}(L_4,C)\xleftarrow{\ \approx\ }\operatorname{Pl}''(L_5,C) & 27\cdot16\cdot10\cdot6=2^6\cdot3^4\cdot5=25\,920\\
\uparrow\ \operatorname{Pl}'(L_5,C) & 27\cdot16\cdot10\cdot6\cdot2=2^7\cdot3^4\cdot5=51\,840\\
\uparrow\ \operatorname{Pl}(L_6,C) & 
\end{array}
\]
\note{colonne de gauche : chaque $\operatorname{Pl}$ est surmonté d'une petite flèche montante vers le précédent ; celle qui précède $\operatorname{Pl}(L_6,C)$ porte « $2$ ». Sous la dernière ligne de nombres, un « $,5$ » isolé. Le $L_1$ de la première ligne est récrit par-dessus une autre lettre.}

On trouve pour les configurations correspondantes
\[
\begin{array}{lll}
B_1(C) & 27:1=27=3^3 & \operatorname{aut}L_1=e\\
B_2(C) & 27\cdot16:2!=27\cdot8=2^3 3^3=216 & \operatorname{aut}L_2=\mathfrak{S}_2\\
B_3(C) & 27\cdot16\cdot10:3!=2^4\cdot3^2\cdot5=720=6! & \operatorname{aut}L_3=\mathfrak{S}_3\\
B_4(C) & 27\cdot16\cdot10\cdot6:4!=2^3\cdot3^3\cdot5=1080 & \operatorname{aut}L_4=\mathfrak{S}_4\\
\quad\simeq B''_4(C) & & \\
B'_5(C) & 27\cdot16\cdot10\cdot6\cdot2:5!=2^4\cdot3^3=432 & \operatorname{aut}L_5=\mathfrak{S}_5\\
B(C)=B_6(C) & 27\cdot16\cdot10\cdot6\cdot2\cdot1:6!=2^3 3^2=72 & \operatorname{aut}L_5\simeq\mathfrak{S}_6\\
\operatorname{Bib}(C) & 36 & \operatorname{aut}(L'_5)\simeq\mathfrak{S}_2\times\mathfrak{S}_6
\end{array}
\]
\note{ainsi sur la page : $B''_4(C)$ (et non $B''_5$), et « $\operatorname{aut}L_5\simeq\mathfrak{S}_6$ » sur la ligne de $B_6$, « $\operatorname{aut}(L'_5)$ » sur celle de $\operatorname{Bib}(C)$ ; la lecture des indices $5$ est sûre, aucun n'est corrigé ici. Dans $2^4\cdot 3^2\cdot 5$ (ligne de $B_3$), le $4$ est écrit sur un $5$. Le signe devant $B''_4(C)$ est un petit « $\simeq$ », de lecture douteuse.}

\emph{Étude \uncertain{cas par cas}} \add{de cardinal 27}

1) $\operatorname{Pl}(L_1,C)=C$ (si $x\in C$, $\operatorname{Aut}(C;x)$ est transitif sur l'ens.\ \add{des 16} \struck{\ill{}} $C-E(x)$ \struck{[\ill{}]})

2) $\operatorname{Pl}(L_2,C)$

\emph{Prop.} Pour \uncertain{tout} couple $(x,y)$ de sommets dist.\ non liés, $\exists!$ base $b$ telle que $x\in b$, $y\in b'$ \add{base opposée} (et $x,y$ \struck{liés} correspondent : $y=\sigma_b x=\sigma_{b'}x$, $x=\sigma_b y=\sigma_{b'}y$). Ainsi l'ens.\ des \add{parties épinglées $2$-libres $(x,y)$} \struck{\ill{}} est en corr.\ 1-1 avec l'ens.\ des couples d'une base $b$ et d'un $x\in b$. La cat.\ des $(C,(x,y))$ équivaut à celle des ens.\ de cardinal $5$ \ill{}, des ens.\ de cardinal $6$ à pt marqué)

On peut écrire $(x,y)=(\xi_6,\xi'_6)$. L'ens.\ $C-E(\xi_6)-E(\xi'_6)-\lbrace x,y\rbrace$ des pts $\neq x,y$ non liés à $x$ ni $y$ est l'ens.\ des $\xi_{ij}$ ($1\leq i<j\leq5$), de cardinal $10=\binom{5}{2}$, et $\operatorname{Aut}(C,(x,y))=\mathfrak{S}_5$ y opère transitivement.

\page{68}
Cas des paires $\lbrace x,y\rbrace$ de pts dist.\ non liés : la donnée équivaut à une \emph{bibase} $B=\lbrace b,b'\rbrace$ avec un élt $z$ dans $B/\sigma_B=I$. La catégorie des $(C,\lbrace x,y\rbrace)$ équivaut à celle des couples $(\tau,J)$ de deux ens.\ avec $\tau$ de cardinal $2$ \add{($\tau=\lbrace b,b'\rbrace$)}, $J$ de cardinal $5$ ($J=B/\sigma_B-\lbrace z\rbrace$).\marginal{$\mathrm{Ens}_2\times\mathrm{Ens}_5$}

\emph{$\operatorname{Pl}(L_3,C)$.} Le type d'une partie \add{épinglée} trois-libre \add{$(x,y,z)$} de $C$ est $(\xi_{12},\xi_{23},\xi_{31})$. L'ens.\ des pts de $C-\lbrace\xi_{12},\xi_{23},\xi_{31}\rbrace$ qui ne sont pas liés à \add{aucun} \struck{aucun} des \uncertain{trois} est $\lbrace\underbrace{\xi_4,\xi_5,\xi_6}\,,\underbrace{\xi'_4,\xi'_5,\xi'_6}\rbrace$\marginal{\ill{}}, qui est de cardinal $6$ ; le groupe des automs de $(C;\xi_{12},\xi_{23},\xi_{31})$ (contenant le groupe des automs de $(C,\lbrace\xi_i\rbrace,\lbrace\xi'_i\rbrace)$ qui respectent la partition \add{de $I=$} $[1,6]$ en $[1,3]$ et $[4,6]$) est transitif sur l'ens.\ précédent.

\emph{Prop.} Pour les couples \add{$(B,K)$} formés d'une bibase $B=\lbrace b,b'\rbrace$ \add{de $C$} et d'une partie $K$ de $B/\sigma_B$ telles que $\operatorname{card}K=3$, soit $J(B,K)$ l'ens.\ partie de $C$ formée des $\xi^B_{ij}$ ($i,j\in K$). On trouve ainsi une bijection de l'ens.\ $\mathcal{B}$ (de cardinal $72\cdot\binom{6}{3}\cdot\frac12=720$) des couples $(B,K)$ avec l'ens.\ des parties $3$-libres de $C$.
\note{sous $\binom{6}{3}$, un « $20$ ».}

La cat.\ des couples $(C,J)$ d'un graphe cubique et d'une partie $3$-libre de $C$ équivaut à celle des triples $(\tau,K,K')$ d'ens.\ de cardinal $2,3,3$ --- on associe à un triple le couple $C=C(\tau,K\amalg K')=(\tau\times(K\amalg K'))\amalg\mathfrak{P}_2(K\amalg K')$, et l'ens.\ $J$ des $\xi^B_{ij}$ ($i,j\in K$), où $B$ est la bibase can.\ de $C$. Le groupe\marginal{$\mathrm{Ens}_2\times\mathrm{Ens}_3\times\mathrm{Ens}_3$}\note{lecture douteuse de la note marginale.}

\page{69}
des automs de $(C,J)$ est donc canoniquement
\[
\mathfrak{S}_\tau\times\mathfrak{S}_K\times\mathfrak{S}_{K'}\quad(\simeq\mathfrak{S}_2\times\mathfrak{S}_3\times\mathfrak{S}_3)
\]
d'ordre $72$.

\emph{NB} Il y a une involution naturelle de la catégorie des $(\tau,K,K')$, savoir
\[
(\tau,K,K')\longmapsto(\tau,K',K).
\]
Elle se traduit en associant à $(C,J)$ le $(C,J')$, où $J'$ est l'ens.\ des élts de $C$ liés à tous les él.\ de $J$. Dans le cas où $J=\lbrace\xi_{12},\xi_{23},\xi_{31}\rbrace$, on a $J'=\lbrace\xi_{45},\xi_{56},\xi_{64}\rbrace$.

\emph{Corollaire.} $J$ admet exactement $2$ bases, \add{savoir $J\cup(b|K')$ et $J\cup(b'|K')$} \struck{(dans le cas envisagé, $\lbrace\xi_{12},\xi_{13},\xi_{34}\rbrace$)} \struck{\ill{}} [dans l'exemple, la réunion de $J=\lbrace\xi_{12},\xi_{23},\xi_{31}\rbrace$ avec $\lbrace\xi_4,\xi_5,\xi_6\rbrace$ ou $\lbrace\xi'_4,\xi'_5,\xi'_6\rbrace$] dont l'intersection se réduit à $J$. \struck{Ainsi}

Considérons l'application $\lbrace\beta',\beta''\rbrace\longmapsto\beta'\cap\beta''$ de l'ens.\ des couples de bases telles que $\operatorname{card}\beta'\cap\beta''=3$, vers l'ens.\ des parties $3$-libres de $C$, \ill{} application est une bijection. \ill{} identifié aussi : l'ens.\ des paires $\lbrace r,s\rbrace$ de racines telles que $\widehat{r,s}=60^\circ$.\marginal{\uncertain{explicitez} \ill{} \uncertain{possibles} \ill{} $\lbrace r,s\rbrace$}

\page{70}
\note{feuille en largeur, écrite sur deux colonnes ; on donne la colonne de gauche, puis celle de droite.}

\emph{Corollaire} (cf.\ plus bas). Pour toute partie $3$-libre $J$ de $C$, l'ens.\ $\mathcal{L}(J)$ des él.\ de $C-J$ non liés aux él.\ de $J$, est un \emph{hexagone}, et on trouve ainsi une corr.\ biun.\ entre $B_3(C)$ et l'ens.\ $\operatorname{Hex}(C)$ des hexagones dans $C$ ; l'application \uncertain{inverse} associe à un hexagone $\mathcal{H}$ l'ens.\ $\mathcal{L}(\mathcal{H})$ des él.\ de $C-\mathcal{H}$ non liés aux él.\ de $\mathcal{H}$.\marginal{$\mathrm{Ens}_2\times\mathrm{Ens}_4$, en regard de ce paragraphe et du suivant, avec une accolade}

\drawing{dans la marge gauche, un hexagone de sommets $\xi'_6$, $\xi_4$, $\xi'_5$, $\xi_6$, $\xi'_4$, $\xi_5$ (en tournant), ses côtés en trait plein, les diagonales intérieures en tirets, avec de petites flèches.}

\emph{NB} La donnée d'un graphe hexagonal (= hexagone combinatoire) équivaut à celle d'un couple \struck{\ill{}} $(\tau,K')$ d'un ens.\ de card.\ $2$ et d'un autre de cardinal $3$ (l'hexagone étant $\tau\times K'$, deux sommets étant liés ssi ils \struck{\ill{}} ont des \struck{\ill{}} \add{composantes} distinctes dans $\tau$ et dans $K'$, i.e.\ ses \struck{\ill{}} parties à deux él.\ de $\tau\times K'$ qui ils forment \uncertain{avec} le graphe d'une \uncertain{bijection} $\tau\hookrightarrow K'$). Ceci dit, si $(C,J)$ est défini par $(\tau,K,K')$, l'hexagone correspondant dans $C$ est défini par $(\tau,K')$, on \ill{}
\[
\tau\times K'\subset\mathcal{B}=\struck{\ill{}}\ \tau\times(K\amalg K')\ \cdots
\]

\emph{Parties $4$-libres.} Cas typique $J=\lbrace\xi_1,\xi_2,\xi_3,\xi_4\rbrace$ ; $\mathcal{L}(J)=\lbrace\xi_5,\xi_6,\xi_{56}\rbrace$. On sait par l'étude précédente que $J$ est contenu dans une base \emph{unique} $b$, et $b\cap\mathcal{L}(J)$ est \uncertain{ici} de cardinal $2$ (\uncertain{c'est} $\lbrace\xi_5,\xi_6\rbrace$), il reste un pt dans le \uncertain{complémentaire} ($\xi_{56}$). Donc $\operatorname{Aut}(C,J)$ (qui est \uncertain{fidèle} \ill{} $\operatorname{Aut}(C,\xi_1,\ldots,\xi_4)$) n'opère pas transitivement sur $\mathcal{L}(J)$, mais \struck{$\mathcal{L}(J)$} a deux orbites $b\cap\mathcal{L}(J)$ ($=\lbrace\xi_5,\xi_6\rbrace$) et le pt fixe $\lbrace\xi_{56}\rbrace$. Donc $\operatorname{Pl}'(L_5,C)\to\operatorname{Pl}(L_4,C)$ est de degré $2$, et $\operatorname{Pl}''(L_5,C)\to\operatorname{Pl}(L_4,C)$ de degré $1$. Notons \ill{} l'ens.\ $E(J)$ des él.\ de $C$ liés \ill{} él.\ de $J$ est de cardinal $2$ (c'est $\lbrace\xi'_5,\xi'_6\rbrace=b'\cap\sigma_b(J)'$ \struck{\ill{}}), et on a une bijection canonique (induite par $\sigma_b$)
\[
E(J)\simeq(b-J)=b\cap\mathcal{L}(J)\subset\mathcal{L}(J).
\]

\emph{Parties $5$-libres.} 1°) Celles qui sont contenues dans une base sont contenues dans une base unique\add{. D'}\struck{\ill{}}\uncertain{Pour} les autres, on vérifie que ce sont des $2$-étoiles.\marginal{dans un cadre, dont une partie est biffée : « Pour des \ill{} de cardinal \ill{} $\mathcal{L}(J)$ \add{de cardinal} \ill{} $E(J)$, \ill{} ($=\lbrace\xi_6\rbrace$), ($=\lbrace\xi'_6\rbrace$), aussi \ill{} $J$ \ill{} La cat.\ des $(C,J)$, $J$ une $5$-base, équivaut à \uncertain{celle} des \struck{couples $(\tau,J)$ d'un ens.\ \ill{} de card.\ $2$} \struck{d'un autre} $J$ de card.\ $5$. » et, contre le cadre, « $\mathrm{Ens}_5$ »}

\emph{Prop.} L'application qui, à toute paire $\lbrace x,y\rbrace$ de pts non liés de $C$ associe $E(x)\cap E(y)$, est une bijection de l'ens.\ des paires avec l'ens.\ des $2$-étoiles (parties $5$-libres ne s'immergeant pas dans une base).\marginal{\add{La bijection inverse est donnée par} $J\mapsto E(J)$} Si on tient compte de $\lbrace x,y\rbrace\leftrightarrow$ (bibase $B=\lbrace b,b'\rbrace$, avec $z\in B/\sigma_B$),

\page{71}
on trouve que $J=\lbrace\xi^B_{z,i}\rbrace_{i\in I-\lbrace z\rbrace}$ --- en termes de base $(\xi_i)_{1\leq i\leq6}$, $z=\lbrace\xi_6,\xi'_6\rbrace$, $J=\lbrace\xi_{i6}\rbrace_{1\leq i\leq5}$. La cat.\ des \struck{\ill{}} couples $(C,J)$, $J$ une $2$-étoile, est donc équivalente à celle des couples $(\tau,J)$, $\tau$ ens.\ de card.\ $2$, $J$ ens.\ de card.\ $5$.\marginal{$\mathrm{Ens}_2\times\mathrm{Ens}_5$}

\emph{Bases et bibases} (il fallait commencer par là !) $W=\operatorname{Aut}(C)$ est simplement transitif sur l'ens.\ des bases épinglées (= ordonnées). Il y a $72$ bases, le groupe de stabilité d'une base est $\simeq\mathfrak{S}_6$. Il y a $36$ bibases, groupe de stabilité $\mathfrak{S}_2\times\mathfrak{S}_6$. La cat.\ des $(C,b)$ est $\simeq$ celle des ens.\ de cardinal $6$, celle des $(C,B=\lbrace b,b'\rbrace)$ à celle des couples $(\tau,I)$, $\tau$ ens.\ de cardinal $2$, $I$ ens.\ de cardinal $6$.\marginal{$\mathrm{Ens}_6$ ; $\mathrm{Ens}_2\times\mathrm{Ens}_6$}

\page{73}
1) \emph{Bitriangle} $\Delta$ : \add{graphe} isomorphe au \add{graphe} produit de deux triangles : $\Delta\simeq S\times T$ \add{($\operatorname{card}S=\operatorname{card}T=3$)}, \add{$x=(\alpha,\beta)$ et $x'=(\alpha',\beta')$, $x\neq y$, liés ssi $\alpha=\alpha'$ ou $\beta=\beta'$ (\uncertain{exclusivement})}, donc $9$ sommets et $4\cdot9:2=18$ arêtes. Les triangles $\subset\Delta$ sont les $6$ fibres des deux projections $S\times T\rightrightarrows S,T$, (\add{\ill{}}) ils se groupent en deux familles par la relation $p\sim q$ ssi $p=q$ ou $p\cap q=\emptyset$ ($\Longleftrightarrow\operatorname{card}p\cap q\neq1$). L'ens.\ des deux familles, $\tau$, donne lieu à un rev.\ d'ordre $3$
\[
T\longrightarrow\tau\qquad(\operatorname{card}\tau=2).
\]
La cat.\ des $\Delta$ équivaut à celle des ens.\ $T$ de cardinal $6$ munis d'une partition de type $(3,3)$, ou encore des flèches de degré $3$ $T\to\tau$ des $\mathrm{Ens}$, avec $\operatorname{card}\tau=2$. On associe à $T\to\tau$ le graphe $\Delta=\prod_{i\in\tau}T_i$, avec \struck{deux arêtes} $x,y$ \add{($\in\Delta$)} liés ssi $x\neq y$ et \struck{\ill{}} $\exists i\in\tau$ tel que $x_i=y_i$ (ou encore, ssi $\exists!\,i\in\tau$, $x_i=y_i$).

2) \emph{Le bloc de triades} (c'est un graphe \struck{\ill{}} $D$ tel que le graphe opposé $\bar D$ ($(x,y)$ liés \add{dans $\bar D$} ssi $x\neq y$ et $x,y$ \emph{non liés} dans $D$) soit un bitriangle. Ainsi $D=S\times T$, $\operatorname{card}S=\operatorname{card}T=3$, avec $x=(\lambda,\mu)$ et $x'=(\lambda',\mu')$ liés ssi $\lambda\neq\lambda'$, $\mu\neq\mu'$. Il y a donc
\[
\operatorname{card}\Bigl(\frac{D\times D-D}{2}\Bigr)-18=18
\]
\struck{arêtes} $18$ arêtes, \ill{} \ill{}, \ill{} que pour le bitriangle $\bar D$. Il y a $6$ triades, i.e.\ \struck{\ill{}} \add{parties} de cardinal $3$ formées de pts mutuellement non liés, qui se groupent en $2$ familles de $3$, d'où un rev.\ de degré $3$ $S\to\sigma$ ($\operatorname{card}\sigma=2$).
\note{la formule de dénombrement est surchargée et en partie biffée ; au-dessus, « $\operatorname{card}D$ » ; lecture douteuse.}

\emph{Proposition.} Le bloc de triades est isomorphe à un bitriangle --- i.e.\ si $\Delta$ \add{($=S\times T$)} est un bitriangle, $\Delta\simeq\bar\Delta$.
\note{l'énoncé est marqué d'un trait vertical dans la marge gauche.}

\emph{Dém.} On va définir dans $\Delta$ une autre famille $\mathfrak{S}$ de $6$ parties à trois éléments --- savoir l'ens.\ $\mathfrak{S}$ des triades $\subset\Delta$. Ce sont aussi les parties $U$ de $\Delta=S\times T$ \add{de cardinal $3$} telles que les deux projections $U\to S$ et $U\to T$ soient injectives (donc bijectives) donc ce sont les graphes des bijections $S\xrightarrow{\sim}T$ --- il y en a bien $6$. Elles se groupent en $2$ familles de trois éléments, en prenant la relation $\varphi\sim\psi$ ssi $\varphi^{-1}\psi$ est une permutation paire de $S$ \struck{condition} [$\mathfrak{S}$ \add{$=\operatorname{Isom}(S,T)$} \struck{\ill{}} torseur sous \struck{$\operatorname{Isom}$ $\operatorname{Aut}$} $\operatorname{Aut}S=\mathfrak{S}_S$, et on prend les $2$ classes sous le sous-groupe d'indice $2$ $\mathfrak{S}^+_S$ ; on voit que ce partage ne dépend pas de l'ordre dans lequel on prend $S,T\ldots$

\page{74}
Si $\varphi,\psi$ \add{($\in\mathfrak{S}_i$ ($i\in\sigma$))} sont dans une même classe \add{($i\in\sigma$)} et $\varphi\neq\psi$, on a $\varphi\cap\psi=\emptyset$, i.e.\ $\varphi^{-1}\psi$ sans pt fixe (en effet, c'est une des deux permut.\ circulaires de $S$) donc les $\varphi\in\mathfrak{S}_i$ forment une \emph{partition} de l'ens.\ $\Delta$ de cardinal $9$. De même pour l'ens.\ $\mathfrak{S}_j$ ($j\in\sigma$, $j\neq i$). On trouve ainsi une application
\[
\Delta\xrightarrow{\ \sim\ }\prod_{i\in\sigma}\mathfrak{S}_i
\]
je dis qu'elle est bijective. Il revient au même de dire qu'elle est injective, ou encore de dire qu'un él.\ $(\alpha,\beta)$ de $\Delta=S\times T$ est connu quand on connaît ses deux composantes dans $\mathfrak{S}_i\times\mathfrak{S}_j$ ($\sigma=\lbrace i,j\rbrace$). Choisissons un isom.\ $S\simeq T$, \add{$\Delta\to S\times S$, $\mathfrak{S}_1$, $\mathfrak{S}_2$} d'où un isom.\ $\sigma\simeq\lbrace1,2\rbrace$, \struck{\ill{}} il revient au même de dire qu'un couple $(\alpha,\beta)$ d'él.\ de $S$ est connu quand on connaît l'unique permutation circulaire $\rho_1$ telle que $\beta=\rho_1\alpha$ (ceci détermine $\beta$ quand $\alpha$ \ill{}) puis l'unique transposition $\theta$ telle que $\beta=(\underbrace{\rho_1\theta}_{\rho_2})\alpha$, i.e.\ telle que $\theta\alpha=\alpha$ --- or $\alpha$ est le \add{(seul) unique} pt fixe de $\theta$.
\note{la glose au-dessus de « Choisissons un isom.\ $S\simeq T$ » est serrée et de lecture douteuse.}

\struck{\ill{}} Il est clair par construction que deux pts de $\Delta$ ayant même image dans \ill{} ou dans \ill{} sont \emph{non liés} dans $\Delta$, donc deux pts de $\Delta$ qui sont liés ont des images distinctes dans \ill{} \emph{et} dans \ill{}, donc sont liés pour la \struck{\ill{}} structure de \emph{bloc de triades} \add{$D$} définie sur $\Delta$ par la bijection $\Delta\simeq\prod_{i\in\sigma}\mathfrak{S}_i$. Comme le cardinal de l'ens.\ des arêtes de $\Delta$ et de $D$ est \uncertain{le même} cardinal $18$, on conclut que les deux structures coïncident. \uncertain{cqfd}.

\emph{Conclusion.} Dans la catégorie \add{(\uncertain{quasi-prédéterminée})} (Bitr) des bitriangles, le foncteur
\[
\Delta\longmapsto\bar\Delta
\]
définit donc une \struck{\ill{}} \add{auto}équivalence involutive. On peut l'interpréter aussi comme une autoéquivalence involutive dans la catégorie des ens.\ \add{$T$} de cardinal $6$ munis d'une \add{partition} \struck{\ill{}} de type $(3,3)$. Pour présenter de façon symétrique cette relation entre $T$ et $S$ de $\mathrm{Ens}_{(3,3)}$, on considère comme « dualité » entre $T$ et $S$\marginal{$\mathrm{Ens}_{\lbrace3,3\rbrace}$}

\page{75}
la donnée d'une \struck{relation} bijection \struck{$R\subset T\times S$} $\prod T\simeq\prod S$ telle que\ldots{} \struck{\ill{} graphe \ill{}}
\note{la ligne biffée est illisible sauf « graphe » ; à droite, deux colonnes de trois points, deux fois. Le reste de la page est blanc.}

\page{77}
\note{page de dénombrements et de figures, sans phrase suivie ; on donne les listes, puis les figures, puis les calculs de la moitié inférieure, écrits tête-bêche.}

polyèdres dans complexe cubique :
\begin{itemize}
\item prismes \quad $45\cdot16$
\item prismes doubles \quad $45\cdot8$
\item antiprismes tordus (appelons \uncertain{projectif} \ill{}) \quad $45\cdot8\cdot16\cdot3=2^4 3^4 5$ \marginal{\ill{} avec \uncertain{graine} \ill{} \uncertain{diagonale} \ill{} \uncertain{Möbius}}
\item pseudocubes (\uncertain{cubes} avec $1$ arête \uncertain{contractée})
\item \emph{cubes} \quad $27\cdot10\cdot2=2^2\cdot3^3\cdot5$
\item dodécaèdre gauche \quad $2^4\cdot3^3=36\cdot12=36\cdot6\cdot2$
\end{itemize}
\note{dans « $45\cdot8\cdot16\cdot3$ » le $8$ est récrit sur un autre chiffre ; dans « $27\cdot10\cdot2$ » le $2$ est écrit sur un $4$. La note en biais, en haut à droite, contre un hexagone divisé par trois diagonales, se lit mal.}

polygones :
\begin{itemize}
\item carrés \quad $27\cdot10\cdot4=2^3 3^3 5$
\item pentagones \quad $36\cdot6\cdot12=2^5\cdot3^4$
\item hexagones \quad $720=2^4\cdot3^2\cdot5$
\end{itemize}

\drawing{en regard des listes : un prisme triangulaire, un prisme double (deux prismes empilés), un « pseudocube » en perspective, et un hexagone traversé de trois diagonales. Au milieu de la page : un dodécaèdre gauche en perspective, ses sommets marqués $\xi_{13}$, $\xi_{45}$, $\xi_{14}$, $\xi_{23}$, $\xi_{51}$, $\xi_{24}$, $\xi_{34}$, $\xi_{12}$, $\xi_{53}$ ; plus loin un pentagone numéroté $1,\ldots,5$ avec la suite « $1\,2\,3\,4\,5$ », un pentagramme dans un carré, une petite figure fermée et un hexagone.}

Sous le dodécaèdre : $\xi_{56}$, $\xi_{26}$, $\xi_{46}$, $\xi_{16}$, $\xi_{36}$ ; $\xi_6$, $\xi'_6$ ; $\xi_{12}$, $\xi_{34}$, $\xi_{51}$, $\xi_{23}$, $\xi_{45}$ (accolade) ; $5\cdot2$ ; puis $2^5\cdot3^4=36\cdot6\cdot12$ pentagones, $2^8\cdot\ldots\cdot5$ \emph{aut.}, et $72=\dfrac{27\cdot16}{6}$.
\note{le $36$ est écrit sur un autre nombre ; « $4!/2=$ » isolé plus à droite.}

Tête-bêche, en bas de la page (barré de deux grands traits obliques) : $C=I\amalg I'\amalg\mathfrak{P}_2(I)$, d'éléments $\xi_i$, $\xi'_i$, $\xi_{ij}$ ; « $27$ » devant $\mathbb{Z}[C]$, une flèche $\mathbb{Z}[C]\to\hat E(C)$, « $7$ » devant $\hat E(C)$, « $6$ » devant $E(C)$, avec la forme
\[
c\cdot c'=\begin{cases}-1 & c=c'\\ 0 & c\neq c',\ c,c'\ \text{non liés}\\ 1 & c\neq c',\ c,c'\ \text{liés}\end{cases}
\]
puis $\xi_1,\ldots,\xi_6$, $\eta=\xi_{ij}+\xi_i+\xi_j$, $\lambda\in\hat E(C)$ avec $\lambda\cdot\sum\xi=1$, et $E(C)=\lambda^{\perp}$.
\note{suivent, du même côté, des ensembles de la forme $\lbrace\xi_i,\ldots\rbrace$ égalés à $9$, une somme $\sum_{b\in\operatorname{Bas}(C)}$ biffée et des bouts de calcul qu'on ne lit pas \ill{}. Lectures de cette moitié douteuses dans l'ensemble.}

\page{78}
\note{page de brouillon, faite surtout de figures ; on relève les formules lisibles, dans l'ordre de haut en bas.}

$M\subset\mathbb{F}_2^{\,Q}$, $P$ ; $Q$ ens.\ de cardinal $5$, $I\subset Q$ ens.\ de cardinal $2$. $\dot e_i\in M$ ; $\dot e_1$, $\dot e_2$, $\dot e_1+\dot e_2$ ;
\[
\dot e_1=e_2+e_3+e_4+e_5,\qquad \dot e_2=e_1+e_3+e_4+e_5,\qquad \dot e_1+\dot e_2=e_1+e_2 ;
\]
$\dot e_i+\mathbb{1}$, $e_j+\mathbb{1}$, $e_i+e_j+\ldots$ ; $e_i+\mathbb{1}$, $e_j+\mathbb{1}$, $e_i+e_j$ ; $\mathbb{F}_2^{\,I}$.
\note{dans $\dot e_1=\ldots$ un terme est surchargé ; le dernier terme de « $e_i+e_j+\ldots$ » est biffé.}

$\tilde I\to I$, $J$ (cardinaux $2$ et $3$) ; \struck{\ill{}} $\tilde I\amalg(J\times\tau(\tilde I/I))$ ; $\underbrace{[0,1]}_{I_0}\amalg I_1\amalg(J\times I_1)$ ; $\dot e_1,\dot e_2,\dot e_3$ \quad \struck{\ill{}} $\mathbb{F}_2^{\,J}$.

\[
3\cdot2^6\qquad 3\cdot2^3\qquad 3\cdot2=6\qquad 2
\]
$1''$ : \uncertain{liés} : $1$ ; non liés : $2'$ (\uncertain{non} : $1'$). $2''$ : liés : $2$ ; liés : $1''$ ; non liés : $1'$, $2'$. \uncertain{liés} $\in\hat E(2)^*$, liés : $1''+\sigma1'$, $\ldots\sigma2'$.
\note{lecture de ce tableau très incertaine.}

\drawing{en haut à gauche, une étoile de dix rayons issus d'un point $s$, les extrémités marquées $\circ$, $\bullet$, $\times$, $*$, $\alpha_1$, $\alpha'_1$, $\alpha_2$, $\alpha'_2$ ; un petit parallélogramme fléché ; un carré aux sommets $\bullet$, $\times$, $\circ$ ; un tableau $2\times2$ de lignes $\alpha_2$, $\alpha'_2$ et colonnes $\alpha_1$, $\alpha'_1$, garni de $\circ$, $*$, $\times$, $\bullet$ ; un cube aux sommets marqués $0$, $e_1$, $e_2$, $e_3$, $e_1+e_2$, $e_1+e_3$, $e_2+e_3$, $e_1+e_2+e_3$ ; un hexagone entouré de trois arcs joignant ses sommets $\alpha_1$, $\alpha'_1$, $\alpha_2$, $\alpha'_2$ à un point $s$ ; deux boîtes en perspective ; une figure d'éventails autour de sommets $1$, $1''$, $2'$, $\alpha'_1$ ; en bas, trois graphes à sommets $u$, $v$, $w$, $\sigma_4$, $\alpha'_3$, $\alpha_4$, $\beta_3$, $\beta_4$, un éventail de triangles, un octaèdre marqué $w$ ; enfin deux lignes $\alpha'_3,\alpha'_4$ / $\beta'_2,\beta'_1,\beta'_3,\beta'_4$ / $\beta_1,\beta_2,\beta_3,\beta'_4$ et, en bas à droite, « $2\cdot6\cdot8$ ».}

\page{79}
\note{page de brouillon, en deux parties séparées par un trait horizontal ; la partie inférieure est barrée de deux traits obliques.}

\[
\underbrace{45\cdot3}\cdot\underbrace{32\cdot3}\ :\ \ldots=45\cdot8\cdot\ldots=2^3\,3^2\,5=360
\]
\note{après le deux-points, un facteur biffé illisible ; le facteur après $45$ est surchargé ; « $2^3$ » est récrit sur « $2^5$ ».}

autom.\ $\ 36\cdot4=144\ \mid\ (\mathfrak{S}_3\times\mathfrak{S}_3)\cdot\mathfrak{S}_2\times\mathfrak{S}_2$.

$W_u(a)$, $w_u(b)$, $w_u(a)-\lbrace b\rbrace$ ; $\underbrace{\xi_{12}\,\xi_{13}\,\xi_{34}}\,\xi_{45}\,\xi_{36}\,\xi_{64}$ ; $\xi_i$, $\xi'_i$, $4\leq i\leq 6$.
\note{la liste des $\xi$ est de lecture douteuse.}

\uncertain{pentaprismes} $\sim$ couples de $3$-bases \uncertain{munis} \struck{\ill{}} \add{d'une} bijection entre les deux --- couples d'une bibase et d'une partition de type $(3,3)$, \uncertain{avec} isom.\ des deux \uncertain{morceaux} i.e.\ isom.\ de $I$ (\uncertain{avec} : $6$ él.) comme produit d'un ens.\ de $3$ par un de $2$.

\drawing{en haut à gauche, un hexagone aux sommets $a$, $b$, $a'$, $b'$ et $u=\infty$, $w=0$ (lecture douteuse), avec une diagonale épaisse ; au-dessous, trois sommets $u$, $v$, $w$ d'où partent des éventails de rayons, marqués $\gamma_4$, $\alpha_1$, $\alpha_2$, $\alpha_3=b'$, $\alpha_4=a$, $\beta_1,\beta'_1,\beta_2,\beta'_2,\beta_3$, $\beta'_3$, $\beta_4=b$ ; un petit prisme étiqueté $1$, $1'$, $2$, $2'$, $3$, $3'$ ; un cube de même ; un schéma $a-b-u$ au-dessus de $a'-b'-u'$, reliés par des traits ; un petit « $\wedge$ ».}

Partie inférieure :
\[
U\mid P_u=\mathrm{id},\quad V\mid P_v=\mathrm{id},\quad W(P_w)=\mathrm{id},\qquad U(\alpha_4)=\alpha_4,\quad V(\beta_4)=\beta_4,\quad W(\gamma_4)=\gamma_4,
\]
\[
UVW\in\lbrace1,\sigma\rbrace\ \overset{?}{\Longrightarrow}\ U,V,W=1 ;
\]
$C\hookrightarrow V\times V\times V$, $N_1$, $N_2$, $N_3$, $C$ ; $J$, $I$ ; $J\simeq V^*$, $I\simeq\lbrace u,v,w\rbrace$ (lecture douteuse) ;
\[
\underset{9}{I\times\lbrace0,1,2\rbrace}\ \amalg\ \underset{3}{I}\times\underset{3}{J}\times\underset{2}{\operatorname{circ}(I)}.
\]
\drawing{à gauche de ces formules, une figure en éventail aux sommets $\gamma_4$, $\beta_4$, $\alpha_4$.}

\page{80}
\note{page de figures, traversée de trois grands traits obliques ; seul le lemme du haut est rédigé.}

\emph{Lemme.} [$u,v,w$ triangle] $u,v$ liés ; $u',u''$ liés à $u$ (et $\neq u,v,w$, i.e.\ non liés à $v$) ; $v',v''$ \ldots\ $v$ (\ldots) ; $u',u''$ liés à chacun $v',v''$ ;
$\lbrace$ $u',v',w'$ triangle ; $u'',v'',w''$ triangle $\rbrace$ $\not\Longrightarrow$
\note{les « \ldots » rendent ses tirets de répétition (« idem ») ; lecture de « chacun » douteuse. La flèche d'implication est barrée d'un trait vertical.}

\drawing{en haut à gauche, un grand polyèdre en perspective aux sommets $u$, $v$, $w$ (triangle en trait épais), $u'$, $u''$, $v'$, $v''$, $w'$, $w''$, $\xi$, $\eta$, $e$, $f$ ; à côté du lemme, deux hexagones, un rectangle surmonté d'un triangle, un « U » fermé par un segment, une chaîne de segments marqués $(u,v)$, $(u',v')$, $(u,v'')$, $(u'',v)$, $(u',v')$, $(u',v)$. Au-dessous, plusieurs cubes et prismes aux sommets $u$, $v$, $u'$, $u''$, $v'$, $v''$, $w'$, $w''$, $\xi$, $\eta$ ; des octaèdres aux sommets $x$, $y$, $s$, $t$, $\xi$, $\eta$ ; un losange aux sommets $\alpha$, $\beta$, $\gamma$ ; un triangle aux sommets $\gamma_4$, $\beta'_4$, $\alpha_4$, $\gamma_1$ ; et, écrites en travers de la page, les listes $\beta_1\,\beta_2\,\beta_3\,\beta_4$, $\beta'_1\,\beta'_2\,\beta'_3\,\beta'_4$, $\gamma_1$, $\gamma_2$, $\gamma_4$, $\alpha_1$, $\alpha'_2$ et « $V\times V\times V$ ».}

\end{document}
